\section{Introduction}

The evolution toward the sixth generation (6G) of wireless networks is characterized by a fundamental paradigm shift from cell-centric architectures to user-centric Cell-Free Massive Multiple-Input Multiple-Output (MIMO) systems. By deploying a large number of distributed Access Points (APs) that cooperatively serve User Equipments (UEs) via a Central Processing Unit (CPU), Cell-Free MIMO offers significant theoretical gains in spectral efficiency and uniform coverage probability. However, the practical realization of these gains is severely constrained by the fronthaul capacity bottleneck. In realistic deployments, the massive volume of in-phase and quadrature (IQ) samples generated by distributed APs must be transmitted to the CPU for joint detection. Since transmitting full-precision IQ samples scales linearly with the number of antennas and bandwidth, it imposes an unsustainable load on limited-capacity fronthaul links. Consequently, the design of efficient quantization and compression schemes that reduce data volume without compromising detection accuracy has become a critical challenge.

Traditionally, fronthaul compression relies on fixed-rate quantization methods, such as uniform quantization or Lloyd-Max algorithms. While these approaches minimize the mean squared error (MSE) between the original and quantized signals, they are fundamentally agnostic to the downstream signal detection task. As noted in \cite{Lewandowsky2021GeneticAT}, the objective of a communication receiver should be to maximize the relevant mutual information rather than simply reconstructing the received waveform. Conventional schemes typically assign equal bit-widths to all APs regardless of their instantaneous Channel State Information (CSI), leading to inefficient resource utilization. For instance, an AP experiencing deep fading or heavy interference contributes little to the final detection but consumes the same fronthaul bandwidth as an AP with a strong line-of-sight connection. Although low-complexity detection algorithms have been analyzed for large-scale systems \cite{Jeon2015OptimalityOL, Som2010ImprovedLD}, they generally assume perfect fronthaul transmission or neglect the non-linear distortion introduced by aggressive quantization.

In recent years, deep learning (DL) has emerged as a powerful tool for physical layer optimization. Model-driven approaches, such as deep unfolding, have been successfully applied to MIMO detection \cite{He2020ModelDriven, Wei2020LearnedCG, Sun2020LearningTS}, demonstrating superior performance over traditional iterative algorithms. However, these works primarily focus on the detection algorithm at the receiver and often overlook the distributed quantization constraints inherent to Cell-Free architectures. To address bandwidth limitations, the concept of Task-Oriented Communication (TOC) has gained traction, often leveraging the Information Bottleneck (IB) principle to extract and transmit only task-relevant features \cite{Shao2021TaskOrientedCF, Shao2021TaskOriented}. For example, \cite{Beck2023SemanticIR} proposes semantic information recovery, while \cite{Li2024TaskOriented} applies Graph Neural Networks (GNNs) and the Graph Information Bottleneck (GIB) to reduce communication overhead in graph data. Similarly, \cite{He2022LearningTO} explores task-oriented channel allocation for multi-agent systems, and \cite{Hou2025Scalable} investigates scalable edge sensing. Despite these advancements, a critical gap remains: existing TOC frameworks often separate the resource allocation decision from the signal detection process or rely on static bit-allocation strategies. There is a lack of a unified framework that allows distributed APs to dynamically infer the optimal quantization precision---including the decision to transmit zero bits---based on local CSI, while the central detector adaptively processes these heterogeneous multi-precision signals.

To bridge this gap, this paper proposes a Joint Quantization-Aware Training (QAT) framework for backhaul-constrained Cell-Free MIMO. Unlike prior works that treat quantization and detection as separate modules, we integrate a lightweight Resource Inference Module at each AP with a Heterogeneous GNN at the CPU. Our approach is grounded in the principle of task-oriented communication: the system is optimized to minimize the final symbol detection error rather than the intermediate quantization distortion. By employing a differentiable bit-selection mechanism, the network learns to allocate higher bit-widths to APs with high-quality observations and silence APs that provide negligible information, thereby maximizing detection performance under strict capacity constraints.

The main contributions of this paper are summarized as follows:
\begin{itemize}
    \item We propose a differentiable bit-selection mechanism utilizing Gumbel-Softmax relaxation, which enables the end-to-end training of discrete quantization levels ($0$, $2$, or $4$ bits). This allows the AP-side Policy Network to dynamically adjust the compression ratio based on instantaneous local CSI.
    \item We design a Bit-Aware GNN detector at the CPU that incorporates quantization precision as an attention bias. This architecture allows the CPU to robustly aggregate signals with heterogeneous quality, effectively compensating for the quantization noise introduced by low-precision APs.
    \item We provide extensive numerical results demonstrating the robustness of the proposed framework. The system maintains a Bit Error Rate (BER) below $2.2\%$ even in extreme scenarios where $50\%$ of APs are disconnected, and achieves detection performance comparable to full-precision baselines while consuming only approximately $2.5$ bits per symbol on average.
\end{itemize}

The remainder of this paper is organized as follows. Section II describes the system model and the problem formulation. Section III details the proposed Joint QAT framework, including the Resource Inference Module and the Bit-Aware GNN. Section IV presents the numerical results and performance analysis. Finally, Section V concludes the paper.